\chapter{Código fuente y cómo ejecutarlo}
\label{anx:codigo}

El código está en el repositorio del equipo. Necesita Python 3.10 o superior. Los algoritmos y
los experimentos solo usan la biblioteca estándar; para las gráficas se necesitan
\texttt{matplotlib} y \texttt{numpy}, que vienen en \texttt{requirements.txt}.

\section{Resolver un grafo}
\begin{verbatim}
python -m src datos/fhcp/graph4.hcp --alg abc --semilla 0
python -m src datos/fhcp/graph4.hcp --alg aco_posa --tour g4.tour
python -m src --listar        # muestra los algoritmos disponibles
\end{verbatim}
El programa carga el grafo, ejecuta la metaheurística y muestra el arreglo de vértices de la
ruta (numerados de 1 a $n$). Si termina con código 0, encontró un ciclo y lo verificó; si termina
con 1, la corrida quedó inconclusa.

\section{Repetir los experimentos y este reporte}
\begin{verbatim}
python -m venv .venv
source .venv/bin/activate          # en fish: source .venv/bin/activate.fish
pip install -r requirements.txt
bash datos/descargar_fhcp.sh       # descarga los 1001 grafos (para B y C)
python -m src.experimentos A --procesos 16     # igual para B y C
python -m src.graficas resultados/disenoA.jsonl
python Reporte/actualizar.py       # pasa tablas y figuras al reporte
cd Reporte && latexmk -pdf main.tex
\end{verbatim}
Si un experimento se interrumpe, basta con volver a lanzar el mismo comando: continúa donde se
quedó.

\section{Archivos principales}
\begin{table}[H]
\centering\small\filas
\begin{tabularx}{\textwidth}{lL}
\toprule
Archivo & Qué contiene \\
\midrule
\texttt{src/grafo.py} & Lectura de grafos HCP y lectura y escritura de rutas TOUR \\
\texttt{src/fitness.py} & Función de costo, cálculo rápido de cambios y verificación del ciclo \\
\texttt{src/representacion.py} & Conversión de \emph{random keys} a ruta (para DE y PSO) \\
\texttt{src/operadores.py} & Inversión guiada por el grafo \\
\texttt{src/algoritmos/} & GA, SA, ACO, ABC, DE y PSO \\
\texttt{src/experimentos.py} & Ejecución de los experimentos A, B y C en paralelo \\
\texttt{src/graficas.py} & Tablas y gráficas a partir de los resultados \\
\texttt{pruebas/} & Pruebas automáticas de cada parte \\
\texttt{Reporte/} & Este documento \\
\bottomrule
\end{tabularx}
\end{table}
